% Generated from locale/tr/content/normal-modal-logic/frame-definability/second-order-definability.tex; do not hand-edit.


\olfileid[tr]{nml}{frd}{st}

\olsection{İkinci Dereceden Tanımlanabilirlik}

Modal formüllerle tanımlanabilen her çerçeve özelliği birinci dereceden
tanımlanabilir değildir. Bununla birlikte bir yerli yüklemler üzerinde
nicelemeye (yani monadik ikinci dereceden nicelemeye) izin verirsek modal
olarak tanımlanabilir bütün çerçeve özelliklerini tanımlayabiliriz. Buradaki
yöntem, modal bir !!{formula}{}ün bir dünyada doğru olduğu koşulları birinci
dereceden !!{formula}s ile ilişkilendirmenin sistemli bir yolundan
yararlanmaktır. Bu, modal formüllerin yalnızca erişilebilirlik bağıntısı için
iki yerli bir $Q$ !!{predicate}{}i değil, bir yerli $P_i$ !!{predicate}{}ini
de içeren bir dildeki birinci dereceden formüllere ``standart çeviri'' adı verilen çevirisidir;
ikinciler $!A$ içinde geçen !!{propositional variable}s $\Obj p_i$ içindir.

\begin{defn}
  \emph{Standart çeviri} $\ST_x(!A)$ tümevarımla şöyle tanımlanır:
  \begin{enumerate}
  \tagitem{prvFalse}{\indcase{!A}{\lfalse}{$\ST_x(\indfrm) = \lfalse$.}}{}
  \tagitem{prvTrue}{\indcase{!A}{\ltrue}{$\ST_x(\indfrm) = \ltrue$.}}{}
  \item \indcase{!A}{\Obj p_i}{$\ST_x(\indfrm) = \Atom{P_i}{x}$.}
  \tagitem{prvNot}{\indcase{!A}{\lnot !B}{$\ST_x(\indfrm) = \lnot \ST_x(!B)$.}}{}
  \tagitem{prvAnd}{\indcase{!A}{(!B \land !C)}{$\ST_x(\indfrm) = (\ST_x(!B)
      \land \ST_x(!C))$.}}{}
  \tagitem{prvOr}{\indcase{!A}{(!B \lor !C)}{$\ST_x(\indfrm) = (\ST_x(!B)
      \lor \ST_x(!C))$.}}{}
  \tagitem{prvIf}{\indcase{!A}{(!B \lif !C)}{$\ST_x(\indfrm) = (\ST_x(!B)
      \lif \ST_x(!C))$.}}{}
  \tagitem{prvIff}{\indcase{!A}{(!B \liff !C)}{$\ST_x(\indfrm) = (\ST_x(!B)
      \liff \ST_x(!C))$.}}{}
  \tagitem{prvBox}{\indcase{!A}{\Box !B}{$\ST_x(\indfrm) =
      \lforall[y][(\Atom{Q}{x,y} \lif \ST_y(!B))]$.}}{}
  \tagitem{prvDiamond}{\indcase{!A}{\Diamond !B}{$\ST_x(\indfrm) =
      \lexists[y][(\Atom{Q}{x,y} \land \ST_y(!B))]$.}}{}
  \end{enumerate}
\end{defn}

Örneğin $\ST_x(\Box p\lif p)$, $\lforall[y][(\Atom{Q}{x,y}
  \lif \Atom{P}{y})]\lif\Atom{P}{x}$ olur. $\ST_x(!A)$'nın diline ait
her !!{structure}, bir evren, $Q$'ya atanan iki yerli bir bağıntı ve bir
yerli $P_i$ !!{predicate}s{}ine atanan evren alt kümeleri gerektirir. Başka
bir deyişle böyle !!a{structure}{}nın bileşenleri, tam olarak $!A$ için bir
modelin bileşenleridir: evren dünyalar kümesidir, $Q$'ya atanan iki yerli
bağıntı erişilebilirlik bağıntısıdır ve $P_i$'ye atanan alt kümeler tam
olarak $V(\Obj p_i)$ atamalarıdır. Modal bir modelde $!A$'nın sağlanmasıyla
karşılık gelen !!{structure} içinde $\ST_x(!A)$'nın sağlanmasının uyuşması
şaşırtıcı değildir:

\begin{prop}\ollabel{prop:st}
  $\mModel{M}=\tuple{W,R,V}$ olsun; $\Struct{M'}$, $\Domain{M'}=W$,
  $\Assign{Q}{M'}=R$ ve $\Assign{P_i}{M'}=V(\Obj p_i)$ olan birinci
  dereceden !!{structure} olsun ve $s(x)=w$ olsun. O zaman
  \[
  \mSat{M}{!A}[w] \text{ ancak ve ancak } \Sat{M'}{\ST_x(!A)}[s]
  \]
\end{prop}

\begin{proof}
  $!A$ üzerinde tümevarımla.
\end{proof}

\begin{prop}
  $!A$'nın modal bir !!{formula} ve
  $\mModel{F}=\tuple{W,R}$'nin bir çerçeve olduğunu varsayalım.
  $\Struct{F'}$, $\Domain{F'}=W$ ve $\Assign{Q}{F'}=R$ olan birinci
  dereceden !!{structure} olsun; ayrıca $!A'$ şu ikinci dereceden formül
  olsun:
  \[
  \lforall[X_1][\dots\lforall[X_n][\lforall[x][
        \SSubst{\ST_x(!A)}{\subst{X_1}{P_1}, \dots,
          \subst{X_n}{P_n}}]]],
  \]
  burada $P_1$, \dots, $P_n$, $\ST_x(!A)$ içindeki bütün bir yerli
  !!{predicate}s{}dir. O zaman
  \[
  \mModel{F} \Entails !A \text{ ancak ve ancak } \Sat{F'}{!A'}
  \]
\end{prop}

\begin{proof}
  $\Sat{F'}{!A'}$ ancak ve ancak $i=1$, \dots, $n$ için
  $\Assign{P_i}{M'}\subseteq W$ olan her !!{structure} $\mModel{M'}$ ve
  $s(x)\in W$ olan her $s$ için $\Sat{M'}{\ST_x(!A)}[s]$ ise geçerlidir.
  \olref{prop:st} gereğince bu, ancak ve ancak $\mModel{F}$'ye dayanan bütün
  $\mModel{M}$ modelleri ve her $w\in W$ dünyası için
  $\mSat{M}{!A}[w]$, yani $\mModel{F}\Entails !A$ ise geçerlidir.
\end{proof}

\begin{defn}
  Bir $\mClass{F}$ çerçeveler sınıfına, ikinci dereceden dilde öyle bir
  !!a{sentence} $!A$ varsa \emph{ikinci dereceden tanımlanabilir} denir:
  bu dil tek bir iki yerli !!{predicate} $P$ ile yalnızca monadik küme
  değişkenleri üzerinde niceleyiciler içerir ve
  $\mModel{F}=\tuple{W,R}\in\mClass{F}$ olması, $\Domain{M}=W$ ve
  $\Assign{P}{M}=R$ olan !!{structure} $\Struct{M}$ içinde
  $\Sat{M}{!A}$ olmasına denktir.
\end{defn}

\begin{cor}
  Bir çerçeveler sınıfı !!a{formula} $!A$ ile tanımlanabiliyorsa, karşılık
  gelen erişilebilirlik bağıntıları sınıfı monadik ikinci dereceden bir
  tümceyle tanımlanabilir.
\end{cor}

\begin{proof}
  Önceki ispatta verilen monadik ikinci dereceden $!A'$ tümcesi gerekli
  özelliğe sahiptir.
\end{proof}

Örnek olarak yine $\Box p\lif p$ !!{formula}{}ü ele alalım. Bu formül
yansımalılığı tanımlar. Yansımalılık elbette $\lforall[x][\Atom{Q}{x,x}]$
!!{sentence}{}yle birinci dereceden tanımlanabilir. Ancak şu monadik ikinci
dereceden !!{sentence} ile de tanımlanabilir:
\[
\lforall[X][\lforall[x][(\lforall[y][(\Atom{Q}{x,y} \lif \Atom{X}{y})]
    \lif \Atom{X}{x})]].
\]
Bu, elbette iki tümcenin denk olduğu anlamına gelir. Bunu doğrudan şöyle
görebilirsiniz: önce ikinci dereceden tümcenin !!a{structure}
$\Struct{M}$ içinde doğru olduğunu varsayalım. $x$ ve $X$ evrensel olarak
nicelendiğinden kalan bölüm, herhangi bir $x\in W$ ve $X\subseteq W$ kümesi,
örneğin $R=\Assign{Q}{M}$ olmak üzere $\Setabs{z}{Rxz}$ kümesi için geçerli
olmalıdır. Dolayısıyla $s(x)\in W$ ve $s(X)=\Setabs{z}{Rxz}$ olan herhangi
bir $s$ için $\Sat{M}{\lforall[y][(\Atom{Q}{x,y}\lif
    \Atom{X}{y})]\lif\Atom{X}{x}}$ olur. Ancak $s(X)$'i seçme biçimimize
göre bu, $\Sat{M}{\lforall[y][(\Atom{Q}{x,y}\lif
    \Atom{Q}{x,y})]\lif\Atom{Q}{x,x}}[s]$ demektir; önbileşen geçerli
olduğundan bu da $\Atom{Q}{x,x}$'e denktir. $s(x)$ gelişigüzel olduğundan
$\Sat{M}{\lforall[x][\Atom{Q}{x,x}]}$ elde edilir.

Şimdi $\Sat{M}{\lforall[x][\Atom{Q}{x,x}]}$ olduğunu varsayalım ve
$\Sat{M}{\lforall[X][\lforall[x][(\lforall[y][(\Atom{Q}{x,y} \lif
        \Atom{X}{y})] \lif \Atom{X}{x})]]}$ olduğunu gösterelim.
Gelişigüzel bir $s$ ataması seçip
$\Sat{M}{\lforall[y][(\Atom{Q}{x,y}\lif\Atom{X}{y})]}[s]$ olduğunu
varsayalım. $s'$, $s'(y)=s(x)$ olan $s$'nin $y$-değişkesi olsun;
$\Sat{M}{\Atom{Q}{x,y}\lif\Atom{X}{y}}[s']$, yani
$\Sat{M}{\Atom{Q}{x,x}\lif\Atom{X}{x}}[s]$ olur.
$\Sat{M}{\lforall[x][\Atom{Q}{x,x}]}$ olduğundan önbileşen doğrudur ve
göstermemiz gereken $\Sat{M}{\Atom{X}{x}}[s]$ sonucunu elde ederiz.

Tanımlanabilir bazı çerçeve sınıfları birinci dereceden tanımlanabilir
olmadığından, $!A'$ biçimindeki her monadik ikinci dereceden !!{sentence}
birinci dereceden bir !!{sentence}{}ye denk değildir. Hangilerinin denk olduğuna
karar veren etkin bir yöntem yoktur.

